\subsection{The Tattvas, the Ages, and the Subtle Body}

The tradition also builds precise numbered systems. The structure of reality, the rhythm of time, the layers of the cosmos, and the makeup of the person are set out as exact enumerations. This part gathers the main ones in tables.

\subsection{The 25 tattvas of Samkhya}

\textbf{Samkhya}, meaning "enumeration," is the classical analytical school. It is dualist. All of reality reduces to two ultimate principles plus twenty-three evolutes of one of them. \textbf{Purusha} is pure consciousness, the witness, inactive and unchanging. It does not create. It only observes. \textbf{Prakriti} is primordial matter or nature, unconscious but active. When its balance is disturbed by the nearness of Purusha, it unfolds the remaining twenty-three principles, called \textbf{tattvas}.

\begin{longtable}{@{} L{0.20\textwidth}L{0.30\textwidth}L{0.40\textwidth}@{}}
\caption{The 25 tattvas of Samkhya, grouped.}\\
\toprule
\textbf{Group} & \textbf{Tattvas} & \textbf{Nature} \\
\midrule
the two ultimates & Purusha, Prakriti & consciousness and primordial matter \\
the inner instrument & Mahat (Buddhi), Ahamkara, Manas & intellect, the I-maker (ego), and the coordinating mind \\
the five knowledge-senses & hearing, touch, sight, taste, smell & the organs of perception \\
the five action-senses & speech, grasping, movement, excretion, generation & the organs of action \\
the five subtle elements & sound, touch, form, taste, smell & the bare sense-objects (tanmatras) \\
the five gross elements & ether, air, fire, water, earth & the physical building blocks (mahabhutas) \\
\bottomrule
\end{longtable}

That is the full set of twenty-five. Liberation in Samkhya is the discriminating knowledge that Purusha is wholly distinct from Prakriti and all its evolutes. Kashmir Shaivism keeps these twenty-five and adds eleven more above them, the pure tattvas and the sheaths of limitation, to make a continuous ladder of thirty-six from pure consciousness down to gross earth.

\subsection{The three gunas}

Prakriti is woven of three strands, the \textbf{gunas}. Their changing proportions make all the variety of nature. In Prakriti's unmanifest state the three are in perfect balance, and creation is the disturbing of that balance.

\begin{longtable}{@{} L{0.18\textwidth}L{0.34\textwidth}L{0.38\textwidth}@{}}
\caption{The three gunas, the strands of nature.}\\
\toprule
\textbf{Guna} & \textbf{Quality} & \textbf{Tendency} \\
\midrule
Sattva & purity, clarity, light & harmony, knowledge, calm \\
Rajas & passion, energy & activity, craving, restlessness \\
Tamas & darkness, heaviness & inertia, ignorance, sleep \\
\bottomrule
\end{longtable}

The Gita holds that even sattva binds. Goodness is still a constraint. The goal is to pass beyond all three.

\subsection{The four yugas and the day of Brahma}

Hindu time is cyclic and immense. The base unit is an age (the \textbf{yuga}). Four yugas make one great age, and dharma, the right order, declines across them in the ratio four to three to two to one.

\begin{longtable}{@{} L{0.24\textwidth}L{0.34\textwidth}L{0.30\textwidth}@{}}
\caption{The four yugas.}\\
\toprule
\textbf{Yuga} & \textbf{Character} & \textbf{Length (human years)} \\
\midrule
Krita (Satya) Yuga & the age of truth, perfect order & 1,728,000 \\
Treta Yuga & order at three quarters & 1,296,000 \\
Dvapara Yuga & order at one half & 864,000 \\
Kali Yuga & the dark age, order at one quarter & 432,000 \\
\bottomrule
\end{longtable}

These four ages nest inside far larger cycles, up to the full life of the creator.

\begin{longtable}{@{} L{0.26\textwidth}L{0.32\textwidth}L{0.30\textwidth}@{}}
\caption{The larger cycles of cosmic time.}\\
\toprule
\textbf{Cycle} & \textbf{Composition} & \textbf{Length} \\
\midrule
Mahayuga & the four yugas together & 4,320,000 years \\
Manvantara & the reign of one Manu & 71 Mahayugas \\
Kalpa, a day of Brahma & 1,000 Mahayugas, 14 Manvantaras & 4.32 billion years \\
a day and night of Brahma & one full day plus its equal night & 8.64 billion years \\
the life of Brahma & 100 years of Brahma & about 311.04 trillion years \\
\bottomrule
\end{longtable}

At the close of each day of Brahma the worlds dissolve, and after his night they form again. At the end of his whole life comes the great dissolution, and all returns to the unmanifest, until a new cycle begins.

\subsection{The 14 lokas}

The cosmos is layered into fourteen worlds, seven upper and seven lower. The seven upper worlds, from the highest down, are \textbf{Satya-loka} (the abode of Brahma), \textbf{Tapa-loka}, \textbf{Jana-loka}, \textbf{Mahar-loka}, \textbf{Svar-loka} (heaven, the world of the gods), \textbf{Bhuvar-loka} (the atmosphere), and \textbf{Bhu-loka} (the earth, the human world). Below the earth lie the seven lower worlds, the patalas, ending in \textbf{Patala} itself, the lowest, abode of the serpents. Below all of these lie the corrective hells.

\subsection{The five koshas and the three bodies}

The person is built in layers. The self is wrapped in five concentric sheaths, the \textbf{koshas}, and these group into three nested \textbf{bodies}.

\begin{longtable}{@{} L{0.26\textwidth}L{0.34\textwidth}L{0.26\textwidth}@{}}
\caption{The five koshas mapped onto the three bodies.}\\
\toprule
\textbf{Kosha (sheath)} & \textbf{Meaning} & \textbf{Body} \\
\midrule
Annamaya & the food sheath, the physical & gross body \\
Pranamaya & the breath or vital sheath & subtle body \\
Manomaya & the mind sheath & subtle body \\
Vijnanamaya & the intellect or wisdom sheath & subtle body \\
Anandamaya & the bliss sheath & causal body \\
\bottomrule
\end{longtable}

The \textbf{gross body} is the physical body of the waking state, made of the five elements, which perishes at death. The \textbf{subtle body} is mind, breath, and intellect together, the body of the dream state, which carries the soul's accumulated tendencies from one life to the next. The \textbf{causal body} is the deepest covering, the seed of ignorance, the body of dreamless sleep, which dissolves only at liberation. Behind all five sheaths and all three bodies stands the \textbf{witness}, the Atman, which does not change through any of the three states. Whatever can be observed is not the observer.

\subsection{The seven chakras and the subtle body}

The Tantric and Hatha Yoga traditions map the subtle body as channels, called \textbf{nadis}, and wheels, called \textbf{chakras}. Three channels are chief. \textbf{Ida} is the cooling lunar channel on the left, \textbf{Pingala} the heating solar channel on the right, and \textbf{Sushumna} the central channel running up the spine. A latent energy called \textbf{kundalini} lies coiled at the base. Through practice it is awakened and rises through Sushumna, piercing each chakra in turn, until it unites with the divine at the crown.

\begin{longtable}{@{} L{0.26\textwidth}L{0.30\textwidth}L{0.18\textwidth}L{0.14\textwidth}@{}}
\caption{The seven chakras of the subtle body.}\\
\toprule
\textbf{Chakra} & \textbf{Location} & \textbf{Element} & \textbf{Seed} \\
\midrule
Muladhara (root) & base of spine & earth & Lam \\
Svadhisthana (sacral) & lower abdomen & water & Vam \\
Manipura (navel) & solar plexus & fire & Ram \\
Anahata (heart) & center of chest & air & Yam \\
Vishuddha (throat) & throat & ether & Ham \\
Ajna (third eye) & between the brows & mind, light & Om \\
Sahasrara (crown) & top of the head & pure consciousness & silent \\
\bottomrule
\end{longtable}

\subsection{Sound as the substance of creation}

One further system runs through all of these. In this tradition sound is not a thing inside the universe. Sound is the substance and origin of the universe. The cosmos is heard before it is seen. The absolute itself is sound, called \textbf{Nada Brahman}. The first sound is the syllable \textbf{Om}, which the tradition splits into its parts, A, U, M, and the silence after, and maps onto waking, dream, deep sleep, and the soundless fourth state beyond all three. The letters of the alphabet are treated as primal sound-powers, the little mothers (the \textbf{matrika}), each one a seed from which a part of reality is born. The petals of the chakras carry these letters, so that the spine holds the whole sound-fabric of the cosmos, and to raise the kundalini is to sound the alphabet from the first letter at the base to the last at the brow and back to silence at the crown.

\subsection{Comparison to the base models}

Set beside the two base models of this collection, Hinduism shows the same pattern stated with unusual detail. There is a single ground beyond being and non-being, a stirring that breaks the one into the many, a layered cosmos and a layered self built in nested shells, and a cyclic rhythm of going out and coming back. Its distinctive features are the breathing of the worlds across immense cycles of time, the precise anatomy of the inner instrument and the sheaths, and the claim that the whole structure is woven of sound. In every register it states the same pattern, the many flowing out from the One and the return of the many to the source.
